%! Polynomial Functions and Rates of Change — Unit 2, Lesson 2.2 — Exit Ticket
\documentclass[10pt]{article}
\usepackage{precalculus-article}
\usepackage{precalculus-boxes}
\newcommand{\TallMath}[1]{$\displaystyle #1\rule[-1.4em]{0pt}{3.2em}$}

\begin{document}

\pageheader{Unit 2, Lesson 2.2}{Exit Ticket}
\namedateperiod

\vspace{0.15in}

\begin{notesbox}{Exit Ticket}
\begin{enumerate}[leftmargin=1.5em, itemsep=10pt]

\item A polynomial function $p$ has exactly three distinct real zeros.

\begin{enumerate}[leftmargin=1.5em, itemsep=6pt]
    \item[(a)] What is the \textbf{minimum number of local extrema} that $p$ must have?

    \vspace{4pt}
    Answer: \blank{3cm}

    \item[(b)] State the class principle that supports your answer.

    \vspace{4pt}
    \blank{11cm}

    \vspace{2pt}
    \blank{11cm}
\end{enumerate}

\item A polynomial function $q$ has a point of inflection at $x = 2$.
The AROC of $q$ over $[0, 2]$ is 3, and the AROC of $q$ over $[2, 4]$ is 7.

\begin{enumerate}[leftmargin=1.5em, itemsep=6pt]
    \item[(a)] Is the AROC of $q$ increasing or decreasing as we move from $[0,2]$ to $[2,4]$?

    \vspace{2pt}
    \blank{9cm}

    \item[(b)] Is the graph of $q$ concave up or concave down on this interval? Explain.

    \vspace{4pt}
    \blank{11cm}

    \vspace{2pt}
    \blank{11cm}
\end{enumerate}

\end{enumerate}
\end{notesbox}

\end{document}
